\subsection{Технология RAG в задачах взаимодействия с базами данных}

\subsubsection{Определение и сущность Retrieval-Augmented Generation (RAG)}

Retrieval-Augmented Generation (RAG), или генерация, дополненная поиском, представляет собой архитектурный подход в области искусственного интеллекта, который объединяет механизмы поиска информации (Retrieval) с генеративными возможностями больших языковых моделей (LLM). Основная концепция RAG заключается в предоставлении LLM доступа к внешним источникам непараметрических знаний, таким как базы данных или корпоративные хранилища документов, непосредственно перед формированием ответа. [RAG-habr]

В рамках этой парадигмы процесс обработки запроса разделяется на два ключевых этапа:

\begin{enumerate}
    \item Поиск (Retrieval): Система RAG извлекает релевантные фрагменты информации (контекст) из внешней базы знаний, используя методы семантического поиска. [энциклопедия]
    \item Генерация (Generation): LLM получает исходный запрос пользователя, дополненный извлеченным контекстом, и на этой основе формирует фактологически точный ответ. [RAG-habr, энциклопедия]
\end{enumerate}

Главная задача RAG - преодоление фундаментальных ограничений LLM, таких как ограниченный объем знаний, зафиксированный на момент обучения, и склонность к "галлюцинациям" (генерации недостоверной информации). Благодаря интеграции актуальных внешних данных RAG значительно повышает фактическую точность и достоверность генерируемых результатов. Кроме того, RAG-системы могут ссылаться на использованные источники, что обеспечивает прозрачность и повышает доверие пользователей. [энциклопедия]

\subsubsection{Эволюция и разновидности архитектур RAG}

Эволюция RAG отражает стремление к созданию более интеллектуальных и ресурсоэффективных систем за счет оптимизации компонентов поиска и генерации.

\textbf{Naive RAG (Наивный RAG)} является начальной методологией, основанной на простом последовательном выполнении трех этапов: индексация, поиск и генерация. Документы преобразуются в векторные эмбеддинги и сохраняются; при запросе ищутся $K$ наиболее релевантных фрагментов по семантическому сходству, которые напрямую передаются LLM. Основной недостаток этой архитектуры — невысокая точность поиска и риск включения нерелевантного контекста, что может приводить к ошибкам генерации.

\textbf{Advanced RAG (Продвинутый RAG)} фокусируется на улучшении качества поиска путем внедрения стратегий оптимизации на этапах до- и после поиска. На этапе pre-retrieval используются такие методы, как оптимизация структуры индекса, уточнение гранулярности данных или переписывание пользовательского запроса (query rewriting) для повышения точности. На этапе post-retrieval применяются техники переранжирования (rerank) и сжатия контекста (context compressing), которые помогают выделить наиболее значимую информацию для LLM.

\textbf{Modular RAG (Модульный RAG)} представляет собой гибкую архитектуру, которая позволяет заменять и реконфигурировать компоненты, отходя от строго последовательного процесса. В эту парадигму вводятся специализированные модули, такие как модуль маршрутизации (Routing), который определяет оптимальный путь обработки запроса, или модуль памяти (Memory) для поддержки многошаговых диалогов.

Особое значение в контексте Text-to-SQL приобретает GraphRAG, который предлагает структурированный и иерархический подход к RAG. Вместо работы с плоскими текстовыми фрагментами GraphRAG извлекает из исходных данных графы знаний (knowledge graphs), организует схему в иерархическую структуру и использует эти отношения для более точного связывания схемы и генерации SQL-запросов. [Welcome graphing] [NL to SQL]

\subsubsection{Проблема интеграции полного контекста базы данных в Text-to-SQL}

Для успешной трансляции запроса на естественном языке (NL) в запрос SQL (Text-to-SQL) LLM должна понимать структуру целевой реляционной базы данных (РБД). [RAG размер и качество] Включение полного контекста РБД, включая полную схему отношений (DDL), описания всех таблиц и столбцов, а также, возможно, примеры данных, в промпт LLM не представляется возможным или оправданным из-за ряда технических и экономических ограничений.

\textbf{3.1. Ограничения контекстного окна и проблема длины промпта}

Одним из ключевых барьеров является лимит контекстного окна (context window) LLM, который определяет максимальное количество токенов, обрабатываемых в одном запросе. В случае сложных корпоративных баз данных, содержащих сотни таблиц с комплексными взаимосвязями, полное DDL-описание и сопутствующие метаданные могут легко превысить этот лимит. [rag-habr]

Даже при наличии расширенных контекстных окон (до 100k+ токенов) передача избыточного объема информации снижает производительность LLM. Если промпт содержит много нерелевантных данных, LLM может "потеряться в середине" (Lost in the middle), что приводит к игнорированию критически важной информации для генерации SQL-запроса, несмотря на её физическое присутствие в промпте. [rag review]

\textbf{3.2. Экономическая неэффективность и стоимость}

Эксплуатационные расходы на использование LLM (особенно проприетарных моделей через API) напрямую коррелируют с количеством обработанных токенов. Передача полного контекста РБД, который может быть объемным, в каждом запросе Text-to-SQL приводит к значительному увеличению потребления токенов и, соответственно, к высокой экономической неэффективности решения. [Text to SQL in era llm, rag-habr] Поскольку LLM приходится обрабатывать большой объем избыточного "шума" в каждом цикле генерации SQL-запроса, стоимость решения становится неоправданно высокой по сравнению с его фактической ценностью. [Text to SQL in era llm, rag-habr]

\textbf{3.3. Безопасность данных и конфиденциальность}

В задачах Text-to-SQL, особенно при работе с корпоративной или конфиденциальной информацией, существует высокий риск нарушения безопасности данных. Если система использует внешние API, отправка полного DDL-описания, а тем более примеров данных (cell values), может привести к нежелательной утечке конфиденциальной информации на внешние серверы [rag размер и качество]. Для соответствия регуляторным требованиям (например, GDPR) и обеспечения конфиденциальности, необходимо минимизировать объем передаваемой информации и гарантировать, что LLM работает только с минимальным, строго необходимым подмножеством метаданных.

\subsubsection{Механизм RAG для извлечения релевантного контекста схемы БД}

Для преодоления проблемы перегрузки контекста RAG-системы в Text-to-SQL используют механизм извлечения только релевантного подмножества схемы БД, который обеспечивает эффективное связывание естественного языка со схемой (Schema Linking). [NL to SQL] [Text to SQL in era llm]

\textbf{Индексация структурных метаданных}

На этапе офлайн-индексации метаданные базы данных, включая DDL-описания таблиц, имена столбцов, описания их атрибутов, взаимосвязи (внешние ключи) и, опционально, примеры данных, фрагментируются и преобразуются в векторные эмбеддинги. [rag-habr] В отличие от традиционного RAG, где индексируются полные документы, здесь индексируются структурные элементы БД. Эти эмбеддинги, отражающие семантическое содержание каждого элемента, сохраняются в векторной базе данных (Vector DB), оптимизированной для поиска ближайших соседей. 

\textbf{Семантический поиск релевантных элементов}

При поступлении NL-запроса пользователя, он также кодируется в векторное представление. Затем механизм поиска выполняет семантический поиск в векторной базе данных, используя метрики близости, такие как косинусное сходство.
Целью этого этапа является извлечение минимального, но достаточного набора элементов схемы, которые релевантны запросу. 

Ключевые извлекаемые элементы:

\begin{itemize}
    \item \textbf{Релевантные таблицы и столбцы:} Система идентифицирует, какие таблицы и столбцы из тысяч доступных наиболее точно соответствуют сущностям и атрибутам, упомянутым в NL-запросе. [Text to SQL in era llm]
    \item \textbf{Метаданные:} Извлекаются описания столбцов, а также примеры значений, что критически важно для генерации точных предикатов в WHERE-клаузах. [Text to SQL in era llm]
    \item \textbf{Связи:} Извлекается информация о связях между релевантными таблицами, что позволяет LLM корректно формулировать операторы JOIN. [Text to SQL in era llm]
\end{itemize}

Таким образом, RAG обеспечивает LLM сфокусированным контекстом, минимизируя шум, снижая потребление токенов и обеспечивая надежную основу для генерации SQL.

\textbf{Генерация SQL-запроса}

На финальном этапе LLM-генератор получает структурированный промпт, включающий исходный NL-запрос, извлеченное RAG-системой подмножество схемы БД и специфические инструкции [rag размер и качество]. Основываясь на этом точном контексте, LLM выполняет задачу Text-to-SQL, генерируя синтаксически корректный SQL-запрос, который максимально соответствует семантике пользовательского намерения [rag размер и качество, nl to SQL].
